% ===========================================================================
%  Homework 02 — Homework 02
%
%  DIVISION OF LABOUR (see AI_INSTRUCTIONS.md §7):
%    The assistant transcribes each assigned problem statement faithfully from
%    the assignment sheet in this folder's assignment/ directory, and emits an
%    empty, marked solution region beneath it.
%    The user types the mathematics inside those regions. No assistant writes
%    inside a SOLUTION region in normal mode.
% ===========================================================================
\documentclass[11pt]{article}
\usepackage{coursemacros}

\title{Homework 02}
\author{Mikey Ferguson}
\date{\today}

\begin{document}
\maketitle

% ---------------------------------------------------------------------------
% Assistant: replace this block with one problem/solution pair per assigned
% problem, in the order assigned, following the pattern below exactly.
% ---------------------------------------------------------------------------

\begin{problem}{1}
  Let $X$, $Y$ and $Z$ be random variables and let $c$ be a number. Prove that
  \[
    \Var(cX) = c^{2}\Var(X),
  \]
  \[
    \Cov(cX, Y) = c\,\Cov(X, Y)
  \]
  and
  \[
    \Cov(X + Y, Z) = \Cov(X, Z) + \Cov(Y, Z).
  \]
\end{problem}
% ===== SOLUTION 1 =====
\paragraph{(i) $\Var(cX) = c^{2}\Var(X)$.}
\begin{align*}
  \Var(cX)
    &= \E\!\left(\left(cX - \E(cX)\right)^{2}\right)\\
    &= \E\!\left(c^{2}X^{2} - 2cX\,\E(cX) + \left(\E(cX)\right)^{2}\right)\\
    &= \E\!\left(c^{2}X^{2} - 2c^{2}X\,\E(X) + \left(c\,\E(X)\right)^{2}\right)\\
    &= \E\!\left(c^{2}X^{2} - 2c^{2}X\,\E(X) + c^{2}\left(\E(X)\right)^{2}\right)\\
    &= \E\!\left(c^{2}\left(X^{2} - 2X\,\E(X) + \left(\E(X)\right)^{2}\right)\right)\\
    &= c^{2}\,\E\!\left(X^{2} - 2X\,\E(X) + \left(\E(X)\right)^{2}\right)\\
    &= c^{2}\,\E\!\left(\left(X - \E(X)\right)^{2}\right)\\
    &= c^{2}\Var(X).
\end{align*}

\paragraph{(ii) $\Cov(cX,Y) = c\,\Cov(X,Y)$.}
Starting from the deviation definition $\Cov(X,Y) = \EE{(X-\mu_{X})(Y-\mu_{Y})}$, and
using $\mu_{cX} = c\,\mu_{X}$:
\begin{align*}
  \Cov(cX, Y)
    &= \EE{(cX - \mu_{cX})(Y - \mu_{Y})}\\
    &= \EE{(cX - c\mu_{X})(Y - \mu_{Y})}\\
    &= \EE{c\,(X - \mu_{X})(Y - \mu_{Y})}\\
    &= c\,\EE{(X - \mu_{X})(Y - \mu_{Y})}\\
    &= c\,\Cov(X,Y).
\end{align*}

\paragraph{(iii) $\Cov(X+Y,Z) = \Cov(X,Z) + \Cov(Y,Z)$.}
\begin{align*}
  \Cov(X+Y, Z)
    &= \EE{\left((X+Y) - \mu_{X+Y}\right)(Z - \mu_{Z})}\\
    &= \EE{(X + Y - \mu_{X} - \mu_{Y})(Z - \mu_{Z})}\\
    &= \EE{\left((X - \mu_{X}) + (Y - \mu_{Y})\right)(Z - \mu_{Z})}\\
    &= \EE{(X - \mu_{X})(Z - \mu_{Z}) + (Y - \mu_{Y})(Z - \mu_{Z})}\\
    &= \EE{(X - \mu_{X})(Z - \mu_{Z})} + \EE{(Y - \mu_{Y})(Z - \mu_{Z})}\\
    &= \Cov(X,Z) + \Cov(Y,Z).
\end{align*}
The collapse $\mu_{X+Y} = \mu_{X} + \mu_{Y}$ on the second line, and the split of the
expectation on the fifth, are both linearity of expectation --- which is exactly what
Problem 2 proves.
% ===== END SOLUTION 1 =====

\begin{problem}{2}
  Let $X_{1}$ and $X_{2}$ be discrete random variables and suppose that $a_{1}$
  and $a_{2}$ are real numbers. Using the joint probability mass function for
  $X_{1}$ and $X_{2}$, show that
  \[
    \E[a_{1}X_{1} + a_{2}X_{2}] = a_{1}\E[X_{1}] + a_{2}\E[X_{2}].
  \]
\end{problem}
% ===== SOLUTION 2 =====
Write $p(x_{1},x_{2})$ for the joint probability mass function of $X_{1}$ and $X_{2}$,
and $p_{X_{1}}$, $p_{X_{2}}$ for the marginals.
\begin{align*}
  \EE{a_{1}X_{1} + a_{2}X_{2}}
    &= \sum_{x_{1}}\sum_{x_{2}} \left(a_{1}x_{1} + a_{2}x_{2}\right) p(x_{1},x_{2})\\
    &= \sum_{x_{1}}\sum_{x_{2}} a_{1}x_{1}\,p(x_{1},x_{2})
       + \sum_{x_{1}}\sum_{x_{2}} a_{2}x_{2}\,p(x_{1},x_{2})\\
    &= \sum_{x_{1}} a_{1}x_{1}\,p_{X_{1}}(x_{1})
       + \sum_{x_{2}}\sum_{x_{1}} a_{2}x_{2}\,p(x_{1},x_{2})\\
    &= \sum_{x_{1}} a_{1}x_{1}\,p_{X_{1}}(x_{1})
       + \sum_{x_{2}} a_{2}x_{2}\,p_{X_{2}}(x_{2})\\
    &= a_{1}\sum_{x_{1}} x_{1}\,p_{X_{1}}(x_{1})
       + a_{2}\sum_{x_{2}} x_{2}\,p_{X_{2}}(x_{2})\\
    &= a_{1}\EE{X_{1}} + a_{2}\EE{X_{2}}.
\end{align*}
The third and fourth lines are marginalization: summing the joint pmf over the variable
that does not appear in the summand leaves the marginal pmf of the other.
% ===== END SOLUTION 2 =====

\begin{problem}{9}
  (Edition 11, Chapter 2.) If the distribution function of $X$ is given by
  \[
    F(b) =
    \begin{cases}
      0            & b < 0,\\
      \frac{1}{2}  & 0 \le b < 1,\\
      \frac{3}{5}  & 1 \le b < 2,\\
      \frac{4}{5}  & 2 \le b < 3,\\
      \frac{9}{10} & 3 \le b < 3.5,\\
      1            & b \ge 3.5,
    \end{cases}
  \]
  calculate the probability mass function of $X$.
\end{problem}
% ===== SOLUTION 9 =====
The mass at each point is the size of the jump of $F$ there, so subtract consecutive
values of $F$: $p(1) = F(1) - F(0)$, and so on.
\[
  p(0) = \tfrac{1}{2},\qquad
  p(1) = \tfrac{1}{10},\qquad
  p(2) = \tfrac{1}{5},\qquad
  p(3) = \tfrac{1}{10},\qquad
  p(3.5) = \tfrac{1}{10},
\]
and $p(x) = 0$ for every other $x$.
% ===== END SOLUTION 9 =====

\begin{problem}{12}
  (Edition 11, Chapter 2.) On a multiple-choice exam with three possible
  answers for each of the five questions, what is the probability that a
  student would get four or more correct answers just by guessing?
\end{problem}
% ===== SOLUTION 12 =====
Five questions, three options each, so the number correct is $\Bin{5}{1/3}$ and
$\PR{\text{correct}} = \tfrac{1}{3}$ on any one question. Four or more correct means
four or five:
\begin{align*}
  \PR{4} + \PR{5}
    &= \binom{5}{4}\left(\tfrac{1}{3}\right)^{4}\left(\tfrac{2}{3}\right)
       + \left(\tfrac{1}{3}\right)^{5}\\
    &= \frac{10}{243} + \frac{1}{243}
     = \boxed{\frac{11}{243}}.
\end{align*}
% ===== END SOLUTION 12 =====

\begin{problem}{22}
  (Edition 11, Chapter 2.) If a fair coin is successively flipped, find the
  probability that a head first appears on the fifth trial.
\end{problem}
% ===== SOLUTION 22 =====
A first head on the fifth trial is exactly the one sequence $TTTTH$, and the flips are
independent and fair, so
\[
  \PR{\text{first head on trial }5} = \left(\tfrac{1}{2}\right)^{5} = \frac{1}{32}.
\]
% ===== END SOLUTION 22 =====

\begin{problem}{37}
  (Edition 11, Chapter 2.) Let $X_1, X_2, \dots, X_n$ be independent random
  variables, each having a uniform distribution over $(0,1)$. Let
  $M = \operatorname{maximum}(X_1, X_2, \dots, X_n)$. Show that the
  distribution function of $M$, $F_M(\cdot)$, is given by
  \[
    F_M(x) = x^n, \qquad 0 \le x \le 1.
  \]
  What is the probability density function of $M$?
\end{problem}
% ===== SOLUTION 37 =====
The maximum is at most $x$ exactly when \emph{every} $X_{i}$ is at most $x$. Concretely,
for $n = 3$ and $x = \tfrac{1}{2}$, the event is $X_{1}\le\tfrac{1}{2}$,
$X_{2}\le\tfrac{1}{2}$, $X_{3}\le\tfrac{1}{2}$, which by independence has probability
$\tfrac{1}{2}\cdot\tfrac{1}{2}\cdot\tfrac{1}{2} = \left(\tfrac{1}{2}\right)^{3}$.
In general, using independence and $\PR{X_{i}\le x} = x$ for $0 \le x \le 1$,
\[
  F_{M}(x) = \PR{M \le x}
           = \PR{X_{1}\le x}\PR{X_{2}\le x}\cdots\PR{X_{n}\le x}
           = x^{n}, \qquad 0 \le x \le 1 .
\]
The curve runs from $0$ to $1$ and is steeper near $1$ as $n$ grows. Differentiating,
\[
  f_{M}(x) = \diff{}{x}\left[x^{n}\right] = n x^{\,n-1}, \qquad 0 \le x \le 1 .
\]
% ===== END SOLUTION 37 =====

% ---------------------------------------------------------------------------
% Remaining assigned work, from the text, Edition 11:
%   problems 61, 76, 86, along with *16 and *49.
%   The starred problems are in the back of the book and are not handed in.
%   Statements for 76, *16, *49 are in the textbook and are
%   not reproduced on the assignment sheet; transcribe each one as it is
%   reached. Problem 86 is printed on the sheet and is transcribed below.
% ---------------------------------------------------------------------------

\begin{problem}{61}
  (Edition 11, Chapter 2.) Let $X$ and $W$ be the working and subsequent
  repair times of a certain machine. Let $Y = X + W$ and suppose that the
  joint probability density of $X$ and $Y$ is
  \[
    f_{X,Y}(x,y) = \lambda^2 e^{-\lambda y}, \qquad 0 < x < y < \infty.
  \]
  \begin{enumerate}
    \item Find the density of $X$.
    \item Find the density of $Y$.
    \item Find the joint density of $X$ and $W$.
    \item Find the density of $W$.
  \end{enumerate}
\end{problem}
% ===== SOLUTION 61 =====
\paragraph{(a) The density of $X$.}
Integrate the joint density over $y$. The region is $0 < x < y < \infty$, so for a fixed
$x$ the variable $y$ runs from $x$ up to $\infty$:
\begin{align*}
  f_{X}(x) &= \int_{y} f_{X,Y}(x,y)\,dy
            = \int_{x}^{\infty} \lambda^{2} e^{-\lambda y}\,dy .
\end{align*}
Substituting $u = -\lambda y$, so that $du = -\lambda\,dy$, $dy = -du/\lambda$, and
$y = x$ becomes $u = -\lambda x$:
\begin{align*}
  f_{X}(x) &= \int_{-\lambda x}^{-\infty} \lambda^{2} e^{u}\left(-\frac{du}{\lambda}\right)
            = \lambda \int_{-\infty}^{-\lambda x} e^{u}\,du
            = \lambda\left[e^{u}\right]_{-\infty}^{-\lambda x}\\
           &= \lambda\left[e^{-\lambda x} - 0\right]
            = \boxed{\lambda e^{-\lambda x}}, \qquad x > 0 .
\end{align*}

\paragraph{(b) The density of $Y$.}
% \todo{label slip on the handwritten page: the left side is written
% $f_Y(x,y)$ where it should be $f_Y(y)$. Mathematics unaffected.}
Integrate the joint density over $x$ instead. For a fixed $y$, the region
$0 < x < y$ makes $x$ run from $0$ to $y$, and $\lambda^{2}e^{-\lambda y}$ is constant
in $x$:
\begin{align*}
  f_{Y}(y) &= \int_{x} f_{X,Y}(x,y)\,dx
            = \int_{0}^{y} \lambda^{2} e^{-\lambda y}\,dx
            = \left[x\,\lambda^{2} e^{-\lambda y}\right]_{0}^{y}\\
           &= \boxed{\lambda^{2} y\, e^{-\lambda y}}, \qquad y \in (0,\infty).
\end{align*}

\paragraph{(c) The joint density of $X$ and $W$.}
Here $Y = X + W$, so substituting $y = x + w$ into $f_{X,Y}$ turns the exponent into a
sum, and the constraint $x < y$ becomes $w > 0$:
\[
  f_{X,W}(x,w) = \boxed{\lambda^{2} e^{-(x+w)\lambda}}, \qquad x > 0,\ w > 0 .
\]

\paragraph{(d) The density of $W$.}
Integrate the joint density of part~(c) over $x$, which now runs from $0$ to $\infty$.
Substituting $u = -(x+w)\lambda$, so $du = -\lambda\,dx$ and $dx = -du/\lambda$:
\begin{align*}
  f_{W}(w) &= \int_{0}^{\infty} \lambda^{2} e^{-(x+w)\lambda}\,dx
            = \int_{-\lambda w}^{-\infty} -\lambda e^{u}\,du
            = \lambda \int_{-\infty}^{-\lambda w} e^{u}\,du\\
           &= \boxed{\lambda e^{-\lambda w}}, \qquad w > 0 .
\end{align*}
\paragraph{(e) $X$ and $W$ are independent.}
Multiplying the two marginals from (a) and (d) reproduces the joint density from (c):
\[
  f_{X}(x)\,f_{W}(w) = \lambda^{2} e^{-\lambda x}\,e^{-\lambda w}
                     = \lambda^{2} e^{-(x+w)\lambda}
                     = f_{X,W}(x,w),
\]
therefore $X$ and $W$ are independent.
% ===== END SOLUTION 61 =====

\begin{problem}{76}
  (Edition 11, Chapter 2.) Let $X$ and $Y$ be independent random variables
  with means $\mu_x$ and $\mu_y$ and variances $\sigma_x^2$ and
  $\sigma_y^2$. Show that
  \[
    \Var(XY) = \sigma_x^2 \sigma_y^2 + \mu_y^2 \sigma_x^2
             + \mu_x^2 \sigma_y^2 .
  \]
\end{problem}
% ===== SOLUTION 76 =====
Start from the computational form of the variance, applied to the single random variable
$XY$, and use independence twice --- once on $\E[X^{2}Y^{2}]$ and once on $\E[XY]$:
\begin{align*}
  \Var[XY] &= \EE{(XY)^{2}} - \left(\EE{XY}\right)^{2}\\
           &= \EE{X^{2}Y^{2}} - \left(\EE{X}\EE{Y}\right)^{2}
              &&\text{($X$, $Y$ independent)}\\
           &= \EE{X^{2}}\EE{Y^{2}} - \E^{2}[X]\cdot\E^{2}[Y].
\end{align*}
Now substitute the second moments, using
$\EE{X^{2}} = \Var(X) + \left(\E(X)\right)^{2} = \sigma_{x}^{2} + \mu_{x}^{2}$ and
likewise for $Y$:
\begin{align*}
  \Var[XY]
    &= \left(\sigma_{x}^{2} + \mu_{x}^{2}\right)
       \left(\sigma_{y}^{2} + \mu_{y}^{2}\right)
       - \mu_{x}^{2}\cdot\mu_{y}^{2}\\
    &= \sigma_{x}^{2}\sigma_{y}^{2} + \mu_{y}^{2}\sigma_{x}^{2}
       + \mu_{x}^{2}\sigma_{y}^{2} + \mu_{x}^{2}\mu_{y}^{2}
       - \mu_{x}^{2}\mu_{y}^{2}\\
    &= \sigma_{x}^{2}\sigma_{y}^{2} + \mu_{y}^{2}\sigma_{x}^{2}
       + \mu_{x}^{2}\sigma_{y}^{2}. \qquad\therefore
\end{align*}
The $+\mu_{x}^{2}\mu_{y}^{2}$ produced by expanding the product cancels the
$-\mu_{x}^{2}\mu_{y}^{2}$ carried down from $\left(\E[XY]\right)^{2}$, leaving the
required three terms.
% ===== END SOLUTION 76 =====

\begin{problem}{86}
  (From Edition 11.) Each new book donated to a library must be processed.
  Suppose that the time it takes a librarian to process a book has mean
  $10$ minutes and standard deviation $3$ minutes. If a librarian has $40$
  books that must be processed one at a time,
  \begin{enumerate}
    \item approximate the probability that it will take more than $420$
      minutes to process all these books;
    \item approximate the probability that at least $25$ books will be
      processed in the first $240$ minutes.
  \end{enumerate}
\end{problem}
% ===== SOLUTION 86 =====
Let $T_{i}$ be the time taken to process book $i$. The $T_{i}$ are independent and
identically distributed with $\EE{T_{i}} = 10$ minutes and $\mathrm{SD}(T_{i}) = 3$
minutes, so $\Var(T_{i}) = 9$ min$^{2}$. Because the books are processed one at a time,
the moment book $n$ is finished is
\[
  S_{n} = T_{1} + T_{2} + \cdots + T_{n},
\]
and by linearity of expectation and independence,
\[
  \EE{S_{n}} = 10n, \qquad \Var(S_{n}) = 9n, \qquad \mathrm{SD}(S_{n}) = 3\sqrt{n}.
\]
The standard deviation is the square root of the variance of the sum; it is not $n$ times
$\mathrm{SD}(T_{i})$. By the central limit theorem, for $n$ this large
\[
  Z = \frac{S_{n} - \EE{S_{n}}}{\mathrm{SD}(S_{n})}
\]
is approximately standard normal, and $\Phi(z) = \PR{Z \le z}$ denotes the standard normal
distribution function.

\paragraph{(a) More than $420$ minutes to process all $40$ books.}
All $40$ books are finished at time $S_{40}$, with
\[
  \EE{S_{40}} = 10(40) = 400 \text{ minutes}, \qquad
  \Var(S_{40}) = 9(40) = 360 \text{ min}^{2}, \qquad
  \mathrm{SD}(S_{40}) = \sqrt{360} = 6\sqrt{10} \approx 18.97 .
\]
Standardizing $420$,
\[
  z = \frac{420 - 400}{6\sqrt{10}} = \frac{20}{6\sqrt{10}} = \frac{10}{3\sqrt{10}}
    = \frac{\sqrt{10}}{3} \approx 1.054 .
\]
The event is ``more than,'' so it is the upper tail:
\begin{align*}
  \PR{S_{40} > 420}
    &\approx \PR{Z > \tfrac{\sqrt{10}}{3}}
     = 1 - \Phi\!\left(\tfrac{\sqrt{10}}{3}\right)\\
    &= 1 - 0.8541
     = \boxed{0.1459}.
\end{align*}

\paragraph{(b) At least $25$ books processed in the first $240$ minutes.}
The normal approximation applies to a sum of processing times, not to a count of books, so
the event must first be rewritten as a statement about $S_{25}$. At least $25$ books are
finished by minute $240$ precisely when the $25$th book is finished by minute $240$, that
is when $S_{25} \le 240$. Note the reversal: \emph{more} books completed corresponds to
\emph{less} time elapsed, so the inequality runs the other way. Here
\[
  \EE{S_{25}} = 10(25) = 250 \text{ minutes}, \qquad
  \Var(S_{25}) = 9(25) = 225 \text{ min}^{2}, \qquad
  \mathrm{SD}(S_{25}) = 15 ,
\]
and standardizing $240$ gives
\[
  z = \frac{240 - 250}{15} = -\frac{10}{15} = -\frac{2}{3} \approx -0.667 .
\]
The event is ``at most,'' so it is the lower tail:
\begin{align*}
  \PR{\text{at least }25\text{ books by minute }240}
    &= \PR{S_{25} \le 240}\\
    &\approx \PR{Z \le -\tfrac{2}{3}}
     = \Phi\!\left(-\tfrac{2}{3}\right)
     = \boxed{0.2525}.
\end{align*}
% ===== END SOLUTION 86 =====

\end{document}
